\documentclass[11pt]{article}
\usepackage[margin=1in]{geometry}
\usepackage{xcolor}
\usepackage[breakable,listings]{tcolorbox}
\tcbuselibrary{listings}
\usepackage{hyperref}
\usepackage{enumitem}
\usepackage{listings}
\usepackage{verbatim}
\usepackage{amsmath}
\usepackage{graphicx}
\usepackage{booktabs}
\usepackage{float}

% Define colors
\definecolor{osaorange}{RGB}{220,115,38}
\definecolor{aiblue}{RGB}{30,144,255}
\definecolor{codebackground}{RGB}{248,248,248}

% Configure hyperref
\hypersetup{
    colorlinks=true,
    linkcolor=blue,
    urlcolor=blue,
    citecolor=blue
}

% Configure listings for code display
\lstset{
  basicstyle=\small\ttfamily,
  breaklines=true,
  breakatwhitespace=true,
  postbreak=\mbox{\textcolor{gray}{$\hookrightarrow$}\space},
  columns=flexible,
  keepspaces=true,
  showstringspaces=false
}

% Code environment using tcolorbox with listings
\newtcblisting{rcode}{
  colback=white,
  colframe=gray!50,
  boxrule=0.5pt,
  arc=0pt,
  left=3pt,
  right=3pt,
  top=3pt,
  bottom=3pt,
  width=\textwidth,
  breakable,
  listing only,
  listing options={
    basicstyle=\small\ttfamily,
    breaklines=true,
    breakatwhitespace=true,
    postbreak=\mbox{\textcolor{gray}{$\hookrightarrow$}\space},
    columns=flexible,
    keepspaces=true,
    showstringspaces=false
  }
}

\title{}
\author{}
\date{}

\begin{document}

% Header box
\begin{tcolorbox}[colback=osaorange,colframe=black,boxrule=2pt,arc=0pt,left=3pt,right=3pt,top=6pt,bottom=6pt,boxsep=2pt]
\centering
\textcolor{white}{\textbf{\Large H524 Introduction to Biostatistics}}\\
\textcolor{white}{\textbf{\Large Week 7 Lab: Nonparametric Methods}}\\
\textcolor{white}{\textbf{\Large Fall 2025}}
\end{tcolorbox}

\vspace{8pt}

\noindent\textbf{Format:} Online\\
\textbf{Tools Required:} R, RStudio\\
\textbf{Estimated Time:} 90 minutes

\section{Learning Objectives}

By the end of this lab, you will be able to:
\begin{enumerate}
\item Check normality assumptions using visual and statistical methods
\item Choose between parametric and nonparametric tests appropriately
\item Perform sign tests in R
\item Use Wilcoxon signed-rank test for one-sample and paired data
\item Conduct Wilcoxon rank-sum tests for independent samples
\item Perform Kruskal-Wallis tests for multiple groups
\item Compare parametric and nonparametric results
\item Interpret rank-based test output
\item \textcolor{aiblue}{\textbf{Use AI tools to help select and implement appropriate tests}}
\end{enumerate}

\section{Introduction: When Assumptions Fail}

Parametric tests (t-tests, ANOVA) assume data are normally distributed. When this assumption is violated---due to skewness, outliers, small sample sizes, or ordinal data---nonparametric tests offer robust alternatives.

\vspace{0.3cm}
\textbf{Key differences:}
\begin{itemize}
\item \textbf{Parametric tests:} Use actual data values, assume normality
\item \textbf{Nonparametric tests:} Use ranks, minimal assumptions
\end{itemize}

\section{Part 1: Checking Normality Assumptions}

Before choosing a test, always check if parametric assumptions hold.

\subsection{Exercise 1.1: Visual Assessment}

\begin{rcode}
# Generate sample data
set.seed(524)
normal_data <- rnorm(30, mean=100, sd=15)
skewed_data <- rexp(30, rate=0.02)  # Exponential (skewed)

# Create multi-panel plot
par(mfrow=c(2,3))

# Normal data
hist(normal_data, main="Normal Data\nHistogram",
     col="lightblue", probability=TRUE)
lines(density(normal_data), col="red", lwd=2)

qqnorm(normal_data, main="Normal Data\nQ-Q Plot")
qqline(normal_data, col="red", lwd=2)

boxplot(normal_data, main="Normal Data\nBoxplot",
        horizontal=TRUE, col="lightblue")

# Skewed data
hist(skewed_data, main="Skewed Data\nHistogram",
     col="lightcoral", probability=TRUE)
lines(density(skewed_data), col="red", lwd=2)

qqnorm(skewed_data, main="Skewed Data\nQ-Q Plot")
qqline(skewed_data, col="red", lwd=2)

boxplot(skewed_data, main="Skewed Data\nBoxplot",
        horizontal=TRUE, col="lightcoral")

par(mfrow=c(1,1))
\end{rcode}

\textbf{Observations:}
\begin{itemize}
\item \textbf{Normal data:} Histogram bell-shaped, Q-Q plot points follow line, boxplot symmetric
\item \textbf{Skewed data:} Histogram right-skewed, Q-Q plot deviates, boxplot asymmetric with outliers
\end{itemize}

\subsection{Exercise 1.2: Shapiro-Wilk Test}

The Shapiro-Wilk test formally tests normality.

\vspace{0.2cm}
\textbf{Hypotheses:}
\begin{itemize}
\item $H_0$: Data are normally distributed
\item $H_A$: Data are not normally distributed
\end{itemize}

\begin{rcode}
# Test normality
shapiro.test(normal_data)
shapiro.test(skewed_data)

# Create function for comprehensive assessment
assess_normality <- function(data, name="Data") {
  cat("\n=== Normality Assessment for", name, "===\n")
  cat("Sample size:", length(data), "\n")
  cat("Mean:", round(mean(data), 2), "\n")
  cat("SD:", round(sd(data), 2), "\n")
  cat("Median:", round(median(data), 2), "\n\n")

  # Shapiro-Wilk test
  sw <- shapiro.test(data)
  cat("Shapiro-Wilk Test:\n")
  cat("  W =", round(sw$statistic, 4), "\n")
  cat("  p-value =", format.pval(sw$p.value, digits=3), "\n\n")

  # Decision
  if (sw$p.value < 0.05) {
    cat("CONCLUSION: Reject normality (p < 0.05)\n")
    cat("RECOMMENDATION: Use nonparametric test\n")
  } else {
    cat("CONCLUSION: Normality assumption reasonable\n")
    cat("RECOMMENDATION: Parametric test appropriate\n")
  }
}

# Use the function
assess_normality(normal_data, "Normal Data")
assess_normality(skewed_data, "Skewed Data")
\end{rcode}

\textbf{Interpretation:}
\begin{itemize}
\item If $p \geq 0.05$: Fail to reject $H_0$, normality assumption reasonable
\item If $p < 0.05$: Reject $H_0$, use nonparametric test
\end{itemize}

\section{Part 2: Sign Test}

The simplest nonparametric test. Uses only the direction of differences.

\subsection{Exercise 2.1: One-Sample Sign Test}

\textbf{Research question:} Is the median recovery time different from 7 days?

\begin{rcode}
# Data: Recovery times (days)
recovery <- c(8, 5, 7, 10, 4, 9, 3, 12, 6, 11, 7)

# Hypothesized median
med0 <- 7

# Manual calculation
diff <- recovery - med0
signs <- sign(diff[diff != 0])  # Remove zeros

n <- length(signs)
b_plus <- sum(signs > 0)
b_minus <- sum(signs < 0)

cat("Number of observations:", length(recovery), "\n")
cat("Non-zero differences:", n, "\n")
cat("Positive differences:", b_plus, "\n")
cat("Negative differences:", b_minus, "\n\n")

# Sign test using binomial test
# H0: median = 7 (p = 0.5 for + or -)
result <- binom.test(b_plus, n, p=0.5, alternative="two.sided")
print(result)

# Interpretation
cat("\nConclusion:\n")
if (result$p.value < 0.05) {
  cat("Reject H0: Median is significantly different from", med0, "\n")
} else {
  cat("Fail to reject H0: No evidence median differs from", med0, "\n")
}
\end{rcode}

\subsection{Exercise 2.2: Paired Sign Test}

\textbf{Before/after treatment comparison}

\begin{rcode}
# Pain scores (1-10 scale)
before <- c(7, 8, 6, 9, 5, 8, 7, 6)
after <- c(5, 6, 4, 7, 4, 6, 5, 5)

# Calculate differences
diff <- before - after
cat("Differences (before - after):", diff, "\n")

# Count signs
signs <- sign(diff)
n_pos <- sum(signs > 0)
n_neg <- sum(signs < 0)
n_zero <- sum(signs == 0)

cat("\nPositive differences:", n_pos, "\n")
cat("Negative differences:", n_neg, "\n")
cat("Zero differences:", n_zero, "\n\n")

# Sign test
result <- binom.test(n_pos, n_pos + n_neg, p=0.5,
                     alternative="two.sided")
print(result)

# Effect size: proportion with improvement
prop_improved <- n_pos / (n_pos + n_neg)
cat("\nProportion improved:", round(prop_improved, 3), "\n")
\end{rcode}

\textbf{Try it yourself:}
Modify the code to test if treatment reduces pain (one-sided test, alternative="greater").

\vspace{2cm}

\section{Part 3: Wilcoxon Signed-Rank Test}

More powerful than sign test---uses both direction AND magnitude of differences.

\subsection{Exercise 3.1: One-Sample Wilcoxon Test}

\begin{rcode}
# Same recovery time data
recovery <- c(8, 5, 7, 10, 4, 9, 3, 12, 6, 11, 7)
med0 <- 7

# Wilcoxon signed-rank test
result <- wilcox.test(recovery, mu=med0, alternative="two.sided")
print(result)

# Summary statistics
cat("\nSummary Statistics:\n")
cat("Sample median:", median(recovery), "\n")
cat("Hypothesized median:", med0, "\n")
cat("Mean:", round(mean(recovery), 2), "\n")
cat("SD:", round(sd(recovery), 2), "\n")

# Compare with sign test
sign_result <- binom.test(sum(sign(recovery - med0) > 0),
                          length(recovery[recovery != med0]),
                          p=0.5, alternative="two.sided")

cat("\n=== Test Comparison ===\n")
cat("Sign test p-value:", format.pval(sign_result$p.value), "\n")
cat("Wilcoxon p-value:", format.pval(result$p.value), "\n")
cat("\nWilcoxon is typically more powerful\n")
\end{rcode}

\subsection{Exercise 3.2: Paired Wilcoxon Test}

\begin{rcode}
# Blood pressure: before and after medication
before <- c(145, 138, 150, 142, 156, 148, 140, 152, 144, 146)
after <- c(138, 135, 142, 140, 148, 145, 136, 145, 141, 143)

# Method 1: Paired test
result1 <- wilcox.test(before, after, paired=TRUE)

# Method 2: Test differences against 0
diff <- before - after
result2 <- wilcox.test(diff, mu=0)

# Both methods equivalent
print(result1)

cat("\nDescriptive Statistics:\n")
cat("Mean difference:", round(mean(diff), 2), "mmHg\n")
cat("Median difference:", median(diff), "mmHg\n")
cat("SD of differences:", round(sd(diff), 2), "\n")

# Visualization
par(mfrow=c(1,2))

boxplot(before, after, names=c("Before", "After"),
        main="Blood Pressure",
        ylab="BP (mmHg)",
        col=c("lightcoral", "lightblue"))

hist(diff, main="Distribution of Differences",
     xlab="Difference (Before - After)",
     col="lightyellow",
     breaks=5)
abline(v=0, col="red", lwd=2, lty=2)
abline(v=median(diff), col="blue", lwd=2)
legend("topright", c("No difference", "Median diff"),
       col=c("red", "blue"), lty=c(2,1), lwd=2)

par(mfrow=c(1,1))
\end{rcode}

\subsection{Exercise 3.3: Compare with Paired t-test}

\begin{rcode}
# Same blood pressure data
before <- c(145, 138, 150, 142, 156, 148, 140, 152, 144, 146)
after <- c(138, 135, 142, 140, 148, 145, 136, 145, 141, 143)
diff <- before - after

# Check normality of differences
shapiro.test(diff)

# Paired t-test
t_result <- t.test(before, after, paired=TRUE)

# Wilcoxon test
w_result <- wilcox.test(before, after, paired=TRUE)

# Compare results
cat("=== Parametric vs. Nonparametric ===\n\n")

cat("Paired t-test:\n")
cat("  t =", round(t_result$statistic, 3), "\n")
cat("  df =", t_result$parameter, "\n")
cat("  p-value =", format.pval(t_result$p.value), "\n")
cat("  95% CI:", round(t_result$conf.int, 2), "\n\n")

cat("Wilcoxon signed-rank test:\n")
cat("  V =", w_result$statistic, "\n")
cat("  p-value =", format.pval(w_result$p.value), "\n\n")

cat("Conclusion: Both tests reach same decision\n")
cat("(when data approximately normal, similar p-values)\n")
\end{rcode}

\section{Part 4: Wilcoxon Rank-Sum Test (Mann-Whitney)}

For comparing two independent groups.

\subsection{Exercise 4.1: Two-Sample Comparison}

\textbf{Clinical trial:} Compare pain relief from two treatments

\begin{rcode}
# Pain relief scores (0-10 scale)
treatment_A <- c(3, 5, 6, 2, 8, 4, 7, 5, 6, 4)
treatment_B <- c(6, 8, 7, 9, 5, 8, 10, 7, 9, 8)

# Check normality
cat("=== Normality Assessment ===\n")
cat("Treatment A:\n")
print(shapiro.test(treatment_A))
cat("\nTreatment B:\n")
print(shapiro.test(treatment_B))

# Wilcoxon rank-sum test
result <- wilcox.test(treatment_A, treatment_B,
                      alternative="two.sided")
print(result)

# Descriptive statistics
cat("\n=== Descriptive Statistics ===\n")
cat("Treatment A:\n")
cat("  Median:", median(treatment_A), "\n")
cat("  Mean:", round(mean(treatment_A), 2), "\n")
cat("  SD:", round(sd(treatment_A), 2), "\n\n")

cat("Treatment B:\n")
cat("  Median:", median(treatment_B), "\n")
cat("  Mean:", round(mean(treatment_B), 2), "\n")
cat("  SD:", round(sd(treatment_B), 2), "\n")

# Visualization
boxplot(treatment_A, treatment_B,
        names=c("Treatment A", "Treatment B"),
        main="Pain Relief Comparison",
        ylab="Pain Relief Score",
        col=c("lightblue", "lightgreen"))
stripchart(list(A=treatment_A, B=treatment_B),
           vertical=TRUE, method="jitter",
           add=TRUE, pch=19, col="darkgray")
\end{rcode}

\subsection{Exercise 4.2: Handling Outliers}

Nonparametric tests are robust to outliers.

\begin{rcode}
# Create data with outlier
set.seed(123)
group_A <- c(23, 25, 28, 30, 32, 35, 38, 95)  # Outlier!
group_B <- c(20, 22, 24, 26, 29, 31, 33, 36)

# Parametric test (affected by outlier)
t_result <- t.test(group_A, group_B)

# Nonparametric test (robust)
w_result <- wilcox.test(group_A, group_B)

# Compare
cat("=== Effect of Outlier ===\n\n")

cat("Descriptive Statistics:\n")
cat("Group A (with outlier=95):\n")
cat("  Mean:", round(mean(group_A), 2), "\n")
cat("  Median:", median(group_A), "\n\n")

cat("Group B:\n")
cat("  Mean:", round(mean(group_B), 2), "\n")
cat("  Median:", median(group_B), "\n\n")

cat("Two-sample t-test:\n")
cat("  p-value =", format.pval(t_result$p.value), "\n")
cat("  Mean difference:", round(diff(t_result$estimate), 2), "\n\n")

cat("Wilcoxon rank-sum test:\n")
cat("  p-value =", format.pval(w_result$p.value), "\n")
cat("  Median A:", median(group_A), "vs Median B:", median(group_B), "\n\n")

cat("Note: Wilcoxon is robust to the extreme outlier (95)\n")
cat("t-test is heavily influenced by it\n")

# Visualization
par(mfrow=c(1,2))

boxplot(group_A, group_B, names=c("A", "B"),
        main="Boxplot (shows outlier)",
        col=c("lightcoral", "lightblue"))
text(1, 95, "Outlier!", pos=4, col="red")

stripchart(list(A=group_A, B=group_B),
           vertical=TRUE, method="jitter",
           main="Individual Points",
           pch=19, col=c("red", "blue"))
points(1, 95, pch=19, col="red", cex=2)

par(mfrow=c(1,1))
\end{rcode}

\section{Part 5: Kruskal-Wallis Test}

Nonparametric alternative to one-way ANOVA.

\subsection{Exercise 5.1: Comparing Three Groups}

\begin{rcode}
# Pain relief from three drugs
drug_A <- c(2, 3, 4, 3, 2, 4)
drug_B <- c(4, 6, 5, 7, 5, 6)
drug_C <- c(5, 8, 9, 10, 7, 9)

# Combine data
pain_relief <- c(drug_A, drug_B, drug_C)
drug <- factor(rep(c("A", "B", "C"), each=6))

# Kruskal-Wallis test
result <- kruskal.test(pain_relief ~ drug)
print(result)

# Descriptive statistics by group
cat("\n=== Descriptive Statistics ===\n")
for (d in c("A", "B", "C")) {
  data <- pain_relief[drug == d]
  cat("\nDrug", d, ":\n")
  cat("  Median:", median(data), "\n")
  cat("  Mean:", round(mean(data), 2), "\n")
  cat("  Range:", range(data)[1], "-", range(data)[2], "\n")
}

# Visualization
boxplot(pain_relief ~ drug,
        main="Pain Relief by Drug",
        xlab="Drug",
        ylab="Pain Relief Score",
        col=c("lightblue", "lightgreen", "lightyellow"))
stripchart(pain_relief ~ drug,
           vertical=TRUE, method="jitter",
           add=TRUE, pch=19, col="darkgray")
\end{rcode}

\subsection{Exercise 5.2: Post-Hoc Tests}

If Kruskal-Wallis is significant, determine which groups differ.

\begin{rcode}
# Same data
pain_relief <- c(drug_A, drug_B, drug_C)
drug <- factor(rep(c("A", "B", "C"), each=6))

# First run Kruskal-Wallis
kw_result <- kruskal.test(pain_relief ~ drug)
cat("Kruskal-Wallis p-value:", format.pval(kw_result$p.value), "\n\n")

if (kw_result$p.value < 0.05) {
  cat("Significant! Now do pairwise comparisons...\n\n")

  # Pairwise Wilcoxon tests with Bonferroni correction
  pairwise_result <- pairwise.wilcox.test(pain_relief, drug,
                                          p.adjust.method="bonferroni")
  print(pairwise_result)

  cat("\nInterpretation:\n")
  cat("- Compare each p-value to alpha = 0.05\n")
  cat("- Bonferroni correction adjusts for multiple tests\n")
} else {
  cat("Not significant. No post-hoc tests needed.\n")
}
\end{rcode}

\subsection{Exercise 5.3: Compare with One-Way ANOVA}

\begin{rcode}
# Same data
pain_relief <- c(drug_A, drug_B, drug_C)
drug <- factor(rep(c("A", "B", "C"), each=6))

# Check normality assumption for ANOVA
cat("=== Check Normality ===\n")
for (d in c("A", "B", "C")) {
  data <- pain_relief[drug == d]
  cat("\nDrug", d, ":\n")
  sw <- shapiro.test(data)
  cat("  Shapiro-Wilk p =", format.pval(sw$p.value), "\n")
}

# ANOVA
anova_result <- aov(pain_relief ~ drug)
cat("\n=== One-Way ANOVA ===\n")
print(summary(anova_result))

# Kruskal-Wallis
kw_result <- kruskal.test(pain_relief ~ drug)
cat("\n=== Kruskal-Wallis Test ===\n")
print(kw_result)

cat("\n=== Comparison ===\n")
cat("ANOVA F-test p-value:",
    format.pval(summary(anova_result)[[1]][["Pr(>F)"]][1]), "\n")
cat("Kruskal-Wallis p-value:", format.pval(kw_result$p.value), "\n")
cat("\nBoth reach same conclusion when assumptions met\n")
\end{rcode}

\section{Part 6: Practical Decision-Making}

\subsection{Exercise 6.1: Test Selection Flowchart}

Create a function that recommends the appropriate test.

\begin{rcode}
recommend_test <- function(data1, data2=NULL, paired=FALSE) {
  cat("=== Test Recommendation System ===\n\n")

  # One-sample or two-sample?
  if (is.null(data2)) {
    cat("Scenario: One-sample test\n\n")

    # Check normality
    sw <- shapiro.test(data1)
    cat("Normality test (Shapiro-Wilk):\n")
    cat("  p-value =", format.pval(sw$p.value), "\n\n")

    if (sw$p.value >= 0.05) {
      cat("RECOMMENDATION: One-sample t-test\n")
      cat("Reason: Data approximately normal\n")
    } else {
      cat("RECOMMENDATION: Wilcoxon signed-rank test\n")
      cat("Reason: Data not normal (p < 0.05)\n")
    }

  } else {
    if (paired) {
      cat("Scenario: Paired two-sample test\n\n")
      diff <- data1 - data2
      sw <- shapiro.test(diff)
      cat("Normality test on differences:\n")
      cat("  p-value =", format.pval(sw$p.value), "\n\n")

      if (sw$p.value >= 0.05) {
        cat("RECOMMENDATION: Paired t-test\n")
        cat("Reason: Differences approximately normal\n")
      } else {
        cat("RECOMMENDATION: Wilcoxon signed-rank test\n")
        cat("Reason: Differences not normal\n")
      }

    } else {
      cat("Scenario: Independent two-sample test\n\n")
      sw1 <- shapiro.test(data1)
      sw2 <- shapiro.test(data2)

      cat("Normality tests:\n")
      cat("  Group 1 p =", format.pval(sw1$p.value), "\n")
      cat("  Group 2 p =", format.pval(sw2$p.value), "\n\n")

      if (sw1$p.value >= 0.05 && sw2$p.value >= 0.05) {
        cat("RECOMMENDATION: Two-sample t-test\n")
        cat("Reason: Both groups approximately normal\n")
      } else {
        cat("RECOMMENDATION: Wilcoxon rank-sum test\n")
        cat("Reason: At least one group not normal\n")
      }
    }
  }
}

# Test the function
set.seed(456)
normal_sample <- rnorm(20, mean=100, sd=15)
skewed_sample <- rexp(20, rate=0.05)

recommend_test(normal_sample)
cat("\n")
recommend_test(skewed_sample)
\end{rcode}

\section{Practice Problems}

\subsection{Problem 1: Hospital Stay Duration}

Compare median hospital stay between two surgical procedures. Data are skewed.

\begin{rcode}
# Hospital stay (days)
procedure_1 <- c(3, 4, 5, 3, 6, 4, 7, 5, 4, 12)
procedure_2 <- c(5, 6, 8, 7, 9, 6, 8, 10, 7, 15)

# Your tasks:
# a) Create boxplots
# b) Check normality
# c) Choose and perform appropriate test
# d) Interpret results in clinical context
\end{rcode}

\textbf{Solution space:}

\vspace{4cm}

\subsection{Problem 2: Ordinal Pain Scores}

Pain scores (1-10) before and after treatment for 15 patients.

\begin{rcode}
# Pain scores (ordinal)
before <- c(7, 8, 6, 9, 5, 8, 7, 6, 8, 7, 9, 6, 7, 8, 5)
after <- c(5, 6, 4, 7, 4, 6, 5, 5, 6, 5, 7, 4, 5, 6, 4)

# Your tasks:
# a) Why is nonparametric test appropriate here?
# b) Perform sign test
# c) Perform Wilcoxon signed-rank test
# d) Compare results and power
\end{rcode}

\textbf{Solution space:}

\vspace{4cm}

\subsection{Problem 3: Multiple Treatment Comparison}

Compare effectiveness of four diet programs.

\begin{rcode}
# Weight loss (kg) after 3 months
diet_A <- c(2.1, 3.2, 2.5, 3.0, 2.8)
diet_B <- c(3.5, 4.2, 3.8, 4.0, 3.9)
diet_C <- c(4.5, 5.1, 4.8, 5.3, 4.9)
diet_D <- c(2.8, 3.0, 2.9, 3.1, 2.7)

# Your tasks:
# a) Combine data appropriately
# b) Perform Kruskal-Wallis test
# c) If significant, do post-hoc tests
# d) Visualize results
# e) Make clinical recommendations
\end{rcode}

\textbf{Solution space:}

\vspace{4cm}

\section{Saving Your Work}

\textbf{IMPORTANT:} Save your R script!

\begin{rcode}
# At the top of your R script, add:
# Name: Your Name
# Date: Today's Date
# Lab: Week 7 - Nonparametric Methods

# Save your workspace
save.image("week7_lab.RData")

# To load it later:
# load("week7_lab.RData")
\end{rcode}

\section{\textcolor{aiblue}{AI-Enhanced Learning Tips}}

\begin{tcolorbox}[colback=aiblue!10,colframe=aiblue,title=Using AI Tools for Nonparametric Tests]

\textbf{Good prompts for AI assistance:}
\begin{itemize}
\item "My data has outliers. Should I use t-test or Wilcoxon? Explain trade-offs"
\item "Show me how to check normality in R with both visual and statistical methods"
\item "Compare parametric and nonparametric results for my data: [paste data]"
\item "Explain when to use sign test vs. Wilcoxon signed-rank test"
\item "Create R code for Kruskal-Wallis with post-hoc tests"
\end{itemize}

\textbf{AI can help with:}
\begin{itemize}
\item Choosing the right test for your data characteristics
\item Generating code for assumption checking
\item Interpreting test output
\item Creating visualizations
\item Understanding when results differ between parametric/nonparametric
\end{itemize}

\textbf{Remember to verify:}
\begin{itemize}
\item Test assumptions AI suggests are actually met
\item P-values and test statistics are reasonable
\item Conclusions match your data characteristics
\item Interpretations make sense in context
\end{itemize}
\end{tcolorbox}

\section{Summary}

\textbf{Key concepts from this lab:}
\begin{enumerate}
\item \textbf{Always check assumptions first}
  \begin{itemize}
  \item Visual: histograms, Q-Q plots, boxplots
  \item Statistical: Shapiro-Wilk test
  \end{itemize}

\item \textbf{Nonparametric test selection}
  \begin{itemize}
  \item One-sample: Sign test or Wilcoxon signed-rank
  \item Paired: Sign test or Wilcoxon signed-rank
  \item Two independent: Wilcoxon rank-sum (Mann-Whitney)
  \item Multiple groups: Kruskal-Wallis
  \end{itemize}

\item \textbf{When to use nonparametric}
  \begin{itemize}
  \item Ordinal data
  \item Non-normal continuous data
  \item Small samples with uncertain normality
  \item Presence of outliers
  \end{itemize}

\item \textbf{Trade-offs}
  \begin{itemize}
  \item Parametric: More powerful IF assumptions met
  \item Nonparametric: Robust, fewer assumptions, slightly less power
  \end{itemize}
\end{enumerate}

\section{Additional Resources}

\begin{itemize}
\item R Documentation:
  \begin{itemize}
  \item \texttt{?wilcox.test} -- Wilcoxon tests
  \item \texttt{?kruskal.test} -- Kruskal-Wallis test
  \item \texttt{?shapiro.test} -- Normality testing
  \end{itemize}
\item Textbook: Pagano \& Gauvreau, Chapter 13
\item Quick-R: Nonparametric Statistics
\end{itemize}

\end{document}
